\chapter{零基础预备知识}

\section{机器人运动学到底在研究什么}

想象桌面上放着一条由六个旋转关节连接起来的机械臂。每个关节都像门轴一样
可以转动。六个关节角写成
\[
\q=[q_1,q_2,q_3,q_4,q_5,q_6].
\]
这里的 $q_i$ 不是坐标点，而是第 $i$ 个电机轴转过的角度。当前项目统一使用
\textbf{弧度}。常见换算为
\[
180^\circ=\pi\ \mathrm{rad},\qquad
1^\circ=\frac{\pi}{180}\ \mathrm{rad}.
\]

机器人学里最常见的两个问题是：

\begin{description}
  \item[正运动学（FK）] 已知六个关节角，计算机械臂末端的位置和朝向。
  \item[逆运动学（IK）] 已知希望末端到达的位置和朝向，反过来求六个关节角。
\end{description}

正运动学像“知道每节手臂怎么弯，算手掌在哪里”；逆运动学像“我想让手掌碰到
杯子，反推肩、肘、腕应该怎么转”。

\section{标量、向量和矩阵}

\subsection{标量}

只有一个数的量称为标量，例如球形障碍物的半径 $r=90\ \mathrm{mm}$。

\subsection{向量}

三个数排成一列可以表示三维空间中的位置：
\[
\p=
\begin{bmatrix}x\\y\\z\end{bmatrix}.
\]
例如 $\p=[480,200,90]^T$ 表示相对于基坐标系，沿 $x$ 方向
\SI{480}{mm}、沿 $y$ 方向 \SI{200}{mm}、沿 $z$ 方向 \SI{90}{mm}。
上标 $T$ 表示转置，即把行向量变成列向量。

向量的长度称为二范数：
\[
\norm{\p}_2=\sqrt{x^2+y^2+z^2}.
\]
代码中的 \code{norm(p)} 就是在算这个量。

\subsection{矩阵}

矩阵可以看成按行列摆放的一组数字。旋转矩阵是一个 $3\times3$ 矩阵：
\[
\R=
\begin{bmatrix}
r_{11}&r_{12}&r_{13}\\
r_{21}&r_{22}&r_{23}\\
r_{31}&r_{32}&r_{33}
\end{bmatrix}.
\]
它的三列分别是新坐标系的 $x$、$y$、$z$ 轴在旧坐标系中指向哪里。

\section{坐标系为什么必不可少}

“末端在 $(100,200,300)$”这句话并不完整，因为必须说明相对于谁测量。
机器人通常给每根连杆建立一个坐标系：
\[
\{0\},\{1\},\ldots,\{6\}.
\]
$\{0\}$ 是固定在机器人底座上的基坐标系，$\{6\}$ 固定在末端。

符号 ${}^0\p_6$ 表示“坐标系 6 的原点，在坐标系 0 中的位置”。
符号 ${}^0\R_6$ 表示“坐标系 6 相对于坐标系 0 的朝向”。

\section{三维旋转}

\subsection{绕三个基本轴旋转}

绕 $x$ 轴旋转 $\alpha$：
\[
\R_x(\alpha)=
\begin{bmatrix}
1&0&0\\
0&\cos\alpha&-\sin\alpha\\
0&\sin\alpha&\cos\alpha
\end{bmatrix}.
\]

绕 $y$ 轴旋转 $\beta$：
\[
\R_y(\beta)=
\begin{bmatrix}
\cos\beta&0&\sin\beta\\
0&1&0\\
-\sin\beta&0&\cos\beta
\end{bmatrix}.
\]

绕 $z$ 轴旋转 $\theta$：
\[
\R_z(\theta)=
\begin{bmatrix}
\cos\theta&-\sin\theta&0\\
\sin\theta&\cos\theta&0\\
0&0&1
\end{bmatrix}.
\]

\begin{warningbox}[矩阵乘法的顺序不能交换]
通常 $\R_x\R_z\ne\R_z\R_x$。先绕 $x$ 转再绕 $z$ 转，与先绕 $z$
再绕 $x$ 转，最终朝向一般不同。因此 MDH 公式中的乘法顺序必须与代码完全一致。
\end{warningbox}

\subsection{旋转矩阵的三个重要性质}

合法旋转矩阵满足
\[
\R^T\R=\boldsymbol{I},\qquad
\R^{-1}=\R^T,\qquad
\det(\R)=1.
\]
第二条非常重要：旋转的逆可以直接用转置得到。项目二中
\[
\R_{36}=\R_{03}^T\R_{06}
\]
正是利用了这一性质。

\section{齐次变换矩阵}

只用 $\R$ 能描述朝向，只用 $\p$ 能描述位置。为了把二者放在一个矩阵里，
机器人学使用 $4\times4$ 齐次变换矩阵：
\[
\T=
\begin{bmatrix}
\R&\p\\
0\ 0\ 0&1
\end{bmatrix}.
\]

若
\[
{}^0\T_1=
\begin{bmatrix}{}^0\R_1&{}^0\p_1\\0&1\end{bmatrix},\quad
{}^1\T_2=
\begin{bmatrix}{}^1\R_2&{}^1\p_2\\0&1\end{bmatrix},
\]
那么从 0 到 2 的变换为
\[
{}^0\T_2={}^0\T_1\,{}^1\T_2.
\]
这就是正运动学循环中不断执行 \code{T = T * A\_i} 的原因。

\section{标准 DH 与改进 DH}

DH 方法用四个参数描述相邻两根连杆：
\[
a,\quad \alpha,\quad d,\quad \theta.
\]
直观解释如下：

\begin{table}[H]
\centering
\begin{tabularx}{0.92\textwidth}{cX}
\toprule
参数 & 含义\\
\midrule
$a$ & 沿某个 $x$ 轴方向的连杆长度\\
$\alpha$ & 相邻两个关节轴之间的扭转角\\
$d$ & 沿关节轴方向的偏移距离\\
$\theta$ & 绕关节轴旋转的角度；转动关节中它含有关节变量\\
\bottomrule
\end{tabularx}
\end{table}

本项目采用改进 DH（MDH）：
\[
{}^{i-1}\T_i=
\R_x(\alpha_{i-1})
\T_x(a_{i-1})
\R_z(\theta_i)
\T_z(d_i).
\]
把四个基本变换相乘，得到代码 \code{mdh\_transform.m} 使用的矩阵：
\[
\begin{bmatrix}
c_\theta&-s_\theta&0&a\\
s_\theta c_\alpha&c_\theta c_\alpha&-s_\alpha&-d s_\alpha\\
s_\theta s_\alpha&c_\theta s_\alpha&c_\alpha&d c_\alpha\\
0&0&0&1
\end{bmatrix}.
\]
其中 $c_\theta=\cos\theta$，$s_\theta=\sin\theta$，其余同理。

\section{关节角偏置 offset}

真实电机零位不一定等于数学模型中 $\theta_i=0$ 的姿态，所以代码使用
\[
\theta_i=q_i+\mathrm{offset}_i.
\]
SR3 的偏置为
\[
\mathrm{offset}=
\left[\pi,\frac{\pi}{2},\frac{\pi}{2},\pi,0,0\right].
\]
例如电机读数 $q_2=0$ 时，MDH 公式实际使用
$\theta_2=\pi/2$。漏掉偏置会让整台机器人模型“从第一步就站错姿势”。

\section{叉乘及其物理意义}

两个三维向量的叉乘为
\[
\boldsymbol{a}\times\boldsymbol{b}=
\begin{bmatrix}
a_yb_z-a_zb_y\\
a_zb_x-a_xb_z\\
a_xb_y-a_yb_x
\end{bmatrix}.
\]
其方向同时垂直于 $\boldsymbol{a}$ 和 $\boldsymbol{b}$。

对转动关节，如果关节轴方向为 $\boldsymbol{z}$，从轴上的一点到末端的向量为
$\p_n-\p_i$，则关节转动产生的瞬时末端线速度方向为
\[
\boldsymbol{z}\times(\p_n-\p_i).
\]
这就是几何雅可比线速度部分的来源。

\section{MATLAB 最低限度语法}

\begin{longtable}{p{0.34\textwidth}p{0.58\textwidth}}
\toprule
代码 & 含义\\
\midrule
\endhead
\code{q(:)'} & 先把数据拉成列向量，再转成行向量，保证形状为 $1\times n$\\
\code{zeros(6,n)} & 创建 $6\times n$ 的全零矩阵\\
\code{eye(4)} & 创建 $4\times4$ 单位矩阵\\
\code{A*B} & 矩阵乘法\\
\code{A.*B} & 对应元素逐个相乘\\
\code{A\textbackslash b} & 解线性方程 $Ax=b$，通常比 \code{inv(A)*b} 稳定\\
\code{T(1:3,4)} & 取齐次矩阵前三行、第四列，即位置向量\\
\code{T(1:3,1:3)} & 取左上角旋转矩阵\\
\code{cross(a,b)} & 三维叉乘\\
\code{norm(x)} & 向量长度；矩阵时需留意所用范数\\
\code{size(A,1)} & 矩阵行数\\
\code{end+1} & 在数组末尾添加新元素\\
\code{struct(...)} & 创建带命名字段的数据结构\\
\bottomrule
\end{longtable}

\begin{plainbox}[读代码的核心方法]
看到一行代码时，先写出每个变量的尺寸。例如
\[
\J:6\times6,\quad \e:6\times1,\quad
\J^T(\J\J^T+\lambda^2I)^{-1}\e:6\times1.
\]
尺寸对得上，代码才可能正确；尺寸对不上，公式一定无法直接运行。
\end{plainbox}
